\chapter*{РЕФЕРАТ}
\addcontentsline{toc}{chapter}{РЕФЕРАТ}

Расчётно-пояснительная записка содержит 60 страниц, 21 иллюстрацию, 7 таблиц, 41 источник, 2 приложения.

Ключевые слова: USB, контроль доступа, доверенные устройства, C++, встраиваемые СУБД, монтирование.

Целью данной дипломной работы является разработка и реализация программно-алгоритмического
комплекса, уменьшающего вероятность несанкционированного копирования информации на сторонние USB-носители.
Результатом работы является клиент-серверное приложение, где серверная часть выполнена на языке C++,
а клиентская на языке TypeScript и фреймворке Vue. Приложение использует встраиваемую СУБД SQlite. Также 
было проведено исследование алгоритма актуализации записей о монтировании при различных сценариях.